\subsection{QP1: técnicas avaliadas e evidências}\label{sec:qp1}
A Tabela~\ref{tab:mecanismos} conta, por tarefa, os estudos em que ao menos uma técnica focal acrescenta cada
mecanismo ao prompt base. Contexto, exemplos, decomposição e restrições são os mecanismos acrescentados com mais
frequência em geração. \textit{Feedback} e raciocínio vêm em seguida. Em depuração, predomina o contexto, que inclui
o código defeituoso, mensagens de erro e informação de localização. Tradução e refatoração têm poucos estudos em
cada célula. A frequência descreve o que os estudos avaliam, e não o quanto cada mecanismo contribui.

\begin{table}[!t]
\caption{Estudos por mecanismo acrescentado pelas técnicas focais}\label{tab:mecanismos}
\centering\footnotesize
\setlength{\tabcolsep}{4pt}
\begin{tabular}{@{}lrrrr@{}}
\hline
Mecanismo & Geração & Depuração & Tradução & Refatoração\\
\hline
Estudos na tarefa & 224 & 59 & 25 & 19\\
\hline
Contexto & 80 & 34 & 4 & 5\\
Exemplos & 72 & 24 & 8 & 8\\
Decomposição & 71 & 18 & 9 & 6\\
Restrições & 68 & 16 & 6 & 3\\
\textit{Feedback} & 66 & 19 & 7 & 5\\
Raciocínio & 59 & 18 & 4 & 4\\
Instrução & 36 & 11 & 1 & 2\\
Multiagente & 26 & 6 & 3 & 1\\
Amostragem & 22 & 8 & 2 & 1\\
Otimização de prompt & 13 & 2 & 2 & 0\\
Outro & 20 & 4 & 2 & 0\\
\hline
\end{tabular}
\end{table}

Das 671 técnicas focais, 349 têm avaliação do pacote, ou seja, a técnica inteira contra a condição sem ela, e 205
são elas mesmas um ingrediente isolado. Outras 3 têm contraste pareado, 13 só têm comparação cumulativa entre
rodadas e 101 não têm condição contrafactual no estudo. As fichas também registram 667 componentes, isto é,
ingredientes das técnicas focais, e 319 deles foram isolados por ablação ou contraste pareado. Por estudo:
\begin{itemize}
  \item 203 dos 304 têm ao menos uma avaliação de pacote;
  \item 181 têm ao menos um ingrediente isolado ou contraste pareado;
  \item 45 não têm condição contrafactual para nenhuma técnica focal ou componente.
\end{itemize}
Entre as técnicas focais, predomina a comparação da técnica completa com uma linha de base, cujo ganho não se
atribui a um mecanismo isolado. As ablações de componentes isolam ingredientes dentro dos pacotes.

As 756 comparações extraídas se distribuem assim, no sentido da métrica de cada uma: 479 com resultado melhor na
condição com a técnica, 138 pior e 26 igual. As outras 94 não têm sentido definido ou não trazem os dois valores, e
19 são descritivas, sem comparador. Esses números não estimam efeito, por quatro razões:
\begin{itemize}
  \item as comparações vêm de \textit{benchmarks}, métricas e modelos diferentes;
  \item a técnica costuma ser a proposta dos próprios autores;
  \item estudos com resultado favorável tendem a ser publicados;
  \item as 42 comparações principais declaradas seguem o recorte dos próprios autores: 38 são favoráveis à técnica,
  e as outras 4 não têm direção definida.
\end{itemize}
A regra de seleção, que escolhe sem olhar o efeito, reduz a escolha seletiva na extração, mas não o viés de
publicação. Por isso, as subseções a seguir dão o mesmo peso aos resultados contrários e às condições em que eles
aparecem.

\subsubsection{Geração de código}

A geração de código é tarefa de 224 estudos. As sínteses parciais cobrem 222 deles, divididos pela técnica focal em três grupos que não se sobrepõem:
\begin{itemize}
\item 92 estudos cuja técnica acrescenta \textit{feedback}, arranjo multiagente ou amostragem de candidatos;
\item 52 com raciocínio ou decomposição, sem esses mecanismos;
\item 78 apenas com instrução, exemplos, contexto, restrições ou otimização do prompt.
\end{itemize}

Os agrupamentos a seguir atravessam os grupos, e um estudo pode aparecer em mais de um. As contagens somam as sínteses parciais de cada grupo; como cada síntese delimitou o agrupamento à sua maneira, as somas são aproximadas e vêm marcadas como ``nas sínteses parciais''. As fichas foram extraídas por IA e ainda não foram conferidas por humano; a qualidade também foi julgada por IA. As comparações foram escolhidas por regra (resultado principal declarado, primeira tabela ou cobertura), e muitas técnicas são propostas dos próprios autores; a proporção de melhoras, portanto, não estima a eficácia geral de uma técnica.

\paragraph{Autorreparo e \textit{feedback} de execução}
Avaliam-se laços em que mensagens do compilador, testes falhos, saídas do interpretador ou alertas de analisadores voltam ao prompt para regenerar o código. Nas sínteses parciais, cerca de 37 estudos entram aqui. A maior parte dos contrastes é de pacote; o laço é isolado como ingrediente numa minoria (entre outros, S150, S479, S645, S878, S1047, S1086, S1898 e S1977), e S917, S1833, S2055 e S2163 são cumulativos.

% AFIRMACAO: A61
Contra um prompt de chamada única, os pacotes de reparo elevaram o Pass@1 de 73,78\% para 90,24\% no HumanEval com gpt-3.5-turbo-0613 ($\mu$FiX; S150-C1) e de 8,73\% para 28,0\% no CoderEval-Python (CoCoGen; S302-C1), e a retirada da fase de \textit{feedback} do $\mu$FiX deixou o valor em 85,37\% (S150-C2).
% AFIRMACAO: A62
Em cinco comparações de ingrediente com valores absolutos, o resultado com o laço foi maior, em métricas que não se agregam: erros de execução de 9 para 6 (S878-C1), HMean \textit{accuracy} de 12,9\% para 18,3\% (S1047-C3), acurácia de 61,7\% para 65,0\% (S1053-C3), Pass@1 de 0,90 para 0,94 sobre modelo ajustado (S1077-C1) e acurácia de 19,6\% para 22,6\% (S1086-C6).
% AFIRMACAO: A63
Nos dois estudos que igualam o orçamento de tentativas ou de programas, não houve ganho: no CodeContests, com pass 1@3, as variantes multiturno com \textit{feedback} ficaram em 2,4\% e 2,8\% contra 2,9\% com Llama 3.0 8B e em 6,9\% contra 7,7\% com Llama 3.1 8B (S479-C2 a S479-C4), e o SEIDR resolveu em média 11,5 problemas do PSB2, contra 12,2 da amostragem sem reparo e 12,0 da busca só de reparo (S645-C1, S645-C2).
% AFIRMACAO: A64
Com Llama 3.2 3B, pôr os testes no prompt rendeu acurácia maior (48,2\%) que o laço de interpretador isolado (22,6\%) e que a combinação dos dois (42,2\%), contra 19,6\% da base (S1086-C5, S1086-C6, S1086-C1).
% AFIRMACAO: A65
Num dos poucos contrastes que isolam a fonte do \textit{feedback}, com 10 requisitos, o relatório estático estruturado e a autorrevisão genérica produziram, cada um, 27 saídas válidas em 30 (S1763-C3).

O laço depende de testes, executor ou compilador disponíveis. Fora de S479 e S645, nenhuma comparação iguala o número de chamadas; uma leitura possível é que parte dos ganhos de pacote acompanhe o acréscimo de chamadas, o que os estudos não isolam. No esclarecimento de requisitos, o usuário simulado vê os testes de referência, condição oráculo (S2140-C2).

\paragraph{Autocrítica e revisão sem execução}
Avalia-se a crítica da saída pelo próprio modelo, ou por relatório externo que não executa o código, seguida de reescrita, muitas vezes dirigida a viés, segurança ou tratamento de exceções. Nas sínteses parciais, são 9 estudos (S114, S140, S149, S397, S604, S759, S834, S842 e S1204), quase todos com contraste de pacote.

% AFIRMACAO: A66
Com PaLM-2-CodeChat-bison, as cinco estratégias de mitigação de viés sem \textit{feedback} elevaram o CBS (menor é melhor) de 17,96\% para valores entre 23,05\% e 33,83\%, e as cinco com relatório de testes de viés o reduziram para valores entre 3,89\% e 6,29\% (S140-C1 a S140-C10).
% AFIRMACAO: A67
A autocrítica RCI reduziu a densidade de fraquezas apontadas pelo Bandit de 0,103 para 0,029 por linha de código com GPT-3 text-davinci-002 (S397-C1) e a proporção de amostras vulneráveis de 5,74\% para 4,5\% com GPT-3.5-turbo (S604-C11), mas, em S397, o teste de Kruskal-Wallis sobre a densidade não foi significativo para GPT-3 (p = 0,334) nem para GPT-3.5 (p = 0,160) (S397-C1).
% AFIRMACAO: A68
Com o modelo pequeno Qwen2.5-1.5B, Self-Refine (38,2\%) e SETS (38,8\%) tiveram Accuracy@t2 menor que a correção simples (44,6\%) (S759-C1, S759-C3), enquanto, em modelos maiores, os ganhos registrados são pequenos, como o CodeBLEU de 0,732 para 0,743 com Gemini-2.0 (S842-C5) e a completude de 92,9\% para 96,0\% com GPT-5 (S1204-C2).

S140 usa um único modelo, e seus autores admitem que parte da redução pode vir só da mudança do prompt. A dependência do porte é hipótese desta revisão, não teste. Analisadores estáticos e juízes LLM não demonstram segurança em uso real.

\paragraph{Sistemas multiagente}
Avaliam-se encadeamentos de agentes com papéis (analista, planejador, testador, revisor) e arranjos em que um modelo externo critica ou implementa. Nas sínteses parciais, são cerca de 29 estudos, todos avaliados como pacote, com ablações de componentes em cerca de dez. Nenhuma comparação com valores contrasta multiagente e agente único com o mesmo número de chamadas.

% AFIRMACAO: A69
Os pacotes elevaram o Pass@1 de 48,1\% para 80,5\% (MapCoder; S344-C1), de 67,68\% para 87,8\% (PairCoder; S448-C1) e de 57,3\% para 74,4\% (Self-collaboration; S2131-C1), e, em S1747, o MapCoder passou de 39,63\% para 64,02\% com Llama3.1 8B usando em média 8,1 chamadas contra 1,0 da base (S1747-C3).
% AFIRMACAO: A70
Em S344, S448, S530 e S634, as oito ablações que retiram um componente registraram valor menor que o do sistema completo, cada uma num único modelo e \textit{benchmark} (S344-C2, S344-C3, S448-C2, S448-C3, S530-C2, S530-C3, S634-C2, S634-C3).
% AFIRMACAO: A71
Há também resultados contrários e nulos: com CodeLlama 34B, o Agents4PLC compilou 31,3\% contra 68,8\% do LLM4PLC (S492-C1); substituir o módulo de codificação progressiva do Spec2RTL-Agent manteve 3 de 3 documentos corretos (S795-C3); e, com Qwen2.5-1.5B, o External-Refine guiado por GPT-4.1 ficou em 41,8\% contra 44,6\% (S759-C2).

Como prompting e orquestração variam juntos, o ganho não se atribui ao mecanismo multiagente. Com modelo externo mais forte (S759-C4, S1204-C3), o efeito se confunde com a seleção de modelo. Há ajuste fino de agentes em S1815, S1828 e S2170, e a única ablação do arranjo multiagente (S530-C3) usa \textit{benchmark} próprio.

\paragraph{Amostragem e seleção de candidatos}
Avalia-se a geração de várias candidatas com escolha por testes, ranqueamento, modelo de valor ou busca em árvore. Nas sínteses parciais, são cerca de 17 estudos, com ingrediente isolado em parte deles (S301, S399, S416, S448, S582, S645, S724 e S1034).

% AFIRMACAO: A72
Em pacotes, o VRank elevou o Pass@1 de 48,8\% para 58,1\% no VerilogEval-Human com GPT-4o-mini (S582-C1) e o SPECFIX, de 87,8\% para 92,0\% no HumanEval+ com DeepSeek-V3 (S724-C1); na ablação deste, retirar a entropia semântica ou a consistência com exemplos levou o valor a 90,55\% e 89,95\%, contra 91,99\% do sistema completo (S724-C2, S724-C3).
% AFIRMACAO: A73
Entre os contrários e nulos, a amostragem com $k = 5$ baixou a criatividade convergente de 2,52 para 2,02 com GPT-3.5 (S399-C3), e a retirada do CoT adaptativo do SEER mudou o pass@1 de 82,0\% para 81,0\%, diferença menor que as dispersões impressas (S1034-C3).

Em S399, a seleção usa o próprio oráculo de avaliação, indisponível sem testes de referência. O SEER inclui modelo de valor treinado, e o TIPCODER, proponente ajustado por SFT e RL; nesses casos (S1034, S1592), o contraste não mede só prompting.

\paragraph{Instrução de raciocínio}
Avalia-se a instrução para raciocinar passo a passo antes do código, ou a inserção de uma cadeia de raciocínio, sobre um prompt que, fora isso, fica igual. Nas sínteses parciais, são cerca de 14 estudos, a maioria com contraste de ingrediente.

% AFIRMACAO: A74
No grupo com raciocínio ou decomposição, as 8 comparações, em 8 estudos, que isolam a instrução de raciocínio como ingrediente dividem-se em 4 melhoras e 4 pioras (S124-C2, S207-C1, S309-C2, S319-C1, S491-C3, S508-C1, S560-C2, S569-C1), e S560 relata ausência de efeito significativo do CoT na correção nos três modelos avaliados (S560-C2).
% AFIRMACAO: A75
No grupo com \textit{feedback}, multiagente ou amostragem, o CoT reduziu o pass@1 de 18,24\% para 14,86\% com Qwen2.5-Coder-7B no BigCodeBench (S1121-C1), de 70,87\% para 61,11\% sobre o prompt EAS com DeepSeek-R1-Distill-Qwen-14B (S1024-C4) e a correção média de 0,53 para 0,49 com Llama-3.1-8B, sob um limite de 512 \textit{tokens} gerados que o estudo associa ao truncamento (S1606-C4), e a elevou de 81,0\% para 87,68\% com DeepSeek-V3 num \textit{benchmark} customizado (S1229-C2).
% AFIRMACAO: A76
Quando a cadeia veio de outro modelo, o pass@1 do Qwen2.5-Coder-7B no BigCodeBench subiu de 18,24\% para 28,38\% com o raciocínio do DeepSeek-R1 (S1121-C4), contraste que não separa o raciocínio do modelo que o produziu.

Comparações de uma célula (um modelo e um \textit{benchmark}) em \textit{benchmarks} distintos não separam o papel do modelo, do domínio e do teto de geração. S207-C1 tem $n = 1$, e S309-C2 mede robustez, não correção.

\paragraph{Planejamento e decomposição}
Avalia-se a geração de plano, pseudocódigo, especificação ou testes intermediários antes do código, num prompt estruturado ou em chamadas encadeadas. Nas sínteses parciais, são cerca de 22 estudos, a maioria com contraste de pacote, além de três estudos de pistas de solução (S892, S1679 e S2145).

% AFIRMACAO: A77
No grupo com raciocínio ou decomposição, 10 contrastes de pacote em 8 estudos registram 8 melhoras e 2 pioras (S1-C1, S6-C1, S77-C1, S216-C1, S216-C2, S295-C2, S491-C1, S654-C1, S654-C2, S845-C1); entre as melhoras estão o Self-planning, de 48,1\% para 60,3\% de Pass@1 no HumanEval com code-davinci-002 contra o prompt direto (S1-C1), e o SCoT, de 53,29\% para 60,64\% contra o CoT (S77-C1).
% AFIRMACAO: A78
As duas pioras ocorreram quando a etapa intermediária veio de um modelo pequeno: a CoT gerada pelo próprio CodeGen 350M baixou o resultado no HumanEval de 14,63\% para 12,8\%, enquanto a de um modelo separado e ajustado o elevou para 20,73\% (S216-C2, S216-C1), e a função auxiliar gerada pelo Incoder 1B mudou o pass@1 de 0,0373 para 0,0364, contra 0,2028 com a função escrita por humanos, condição oráculo (S295-C2, S295-C3).
% AFIRMACAO: A79
Decompor também piorou resultados: com WizardCoder-15B, a geração holística de uma classe Java teve Completion@1 de 98,0\%, contra 74,0\% na síntese independente por método e 70,0\% na incremental, métrica que não mede correção (S365-C1, S365-C3), e o Knowledge Generation Prompting reduziu o pass@1 do GPT-3.5 Turbo no ChemEval de 45,0\% para 39,79\%, em contradição com a narrativa do estudo (S2075-C2).
% AFIRMACAO: A80
Nos dois estudos de pistas de solução com números, os maiores ganhos vieram de informação externa ao modelo avaliado: com DeepSeek-Coder-V2 no MuPA, as pistas autogeradas levaram o pass@1 (escala 0--100) de 11,84 para 13,21, 12,18 e 11,92, e a solução oficial fornecida, para 38,79 (S2145-C1, S2145-C3, S2145-C5, S2145-C6); com Deepseek-coder 33B, os casos de teste aprovados passaram de 269 para 299 com as dicas oficiais e para 801 com dicas reescritas pelo GPT-4o (S1679-C1, S1679-C2).

Os ganhos de pacote não se atribuem à decomposição, e nenhuma comparação mede custo, embora várias técnicas usem mais de uma chamada. Os valores de ``até 25,4\%'' de S1 e ``até 13,79\%'' de S77 são variações relativas máximas, não p.p. As pistas de referência medem o aproveitamento de conhecimento fornecido, não a capacidade de produzi-lo.

\paragraph{Exemplos e sua seleção}
Avaliam-se a presença, a quantidade, a seleção e o formato de demonstrações ou de casos de teste no prompt. Nas sínteses parciais, cerca de 14 estudos do grupo só com instrução, exemplos, contexto ou restrições têm exemplos como técnica focal. Pacotes com seleção treinada ou exploração simbólica (S170-C1, S195-C1, S500-C1) incluem componentes fora do escopo de prompting.

% AFIRMACAO: A81
Quando apenas se acrescentam exemplos, sem recuperação treinada, as diferenças registradas são pequenas: Pass@1 de 0,229 para 0,237 (proporção) no VHDL-Eval (S361-C1), de 59,0\% para 62,8\% no VerilogEval com GPT-4o (S432-C1) e, em R com DeepSeek Coder 1B, de 13,9\% para 14,1\% com \textit{few-shot} e para 13,8\% com exemplos de tradução (S591-C3, S591-C1).
% AFIRMACAO: A82
Os exemplos também pioraram resultados: exemplos de CVT reduziram o pass@1 funcional de 80,81\% para 70,87\% (S1024-C2), amostras de código com erros o derrubaram de 45,0\% para 15,62\% com GPT-3.5 Turbo (S2075-C1), e o \textit{few-shot} elevou a taxa de inacessibilidade do GPT-4o de 0,326 para 0,469 (S2113-C2).

O sinal varia com modelo, temperatura e tarefa (S432, S1707, S2090). Em S902-C2, o conteúdo recuperado recebeu nomes extraídos da solução de referência, condição oráculo, e em S1613-C1 metade dos testes de avaliação foi exposta no prompt.

\paragraph{Contexto recuperado e de repositório}
Avalia-se o acréscimo de documentação de API, código do arquivo ou do repositório, dependências, tipos, dados ou regras de domínio. Nas sínteses parciais, são cerca de 17 estudos do grupo só com instrução, exemplos, contexto ou restrições e 8 do grupo com raciocínio ou decomposição que combinam recuperação e decomposição. Predominam contrastes de pacote; as ablações de S360 e S852 isolam componentes dentro de sistemas completos.

% AFIRMACAO: A83
Com contexto de arquivo ou de repositório, o Pass@1 passou de 17,4\% para 53,04\% no DevEval com gpt-4 (S357-C1) e de 20,29\% para 44,93\% com o RepoScope, que usou em média 3679 \textit{tokens} de prompt contra 97 da base (S852-C1).
% AFIRMACAO: A84
No grupo com raciocínio ou decomposição, 8 das 17 comparações de contexto isolado usam informação oráculo, como o contexto do arquivo somado às APIs invocadas, que elevou o Pass@1 do GPT-4o Mini no RepoExec de 16,62\% para 37,75\% (S213-C2, S213-C3, S213-C5, S295-C3, S670-C2, S670-C5, S670-C6, S670-C7).
% AFIRMACAO: A85
Há inversões e nulos: a recuperação reduziu a acurácia de Claude 3.7 no QHackBench de 31,0\% para 0,0\% (S806-C1), o RAG sobre a documentação do Ansible baixou o BLEU do GPT-4 de 0,1828 para 0,1568 (S1951-C1), a QA API isolada reduziu a acurácia de execução de 59,6\% para 57,4\% (S176-C3) e o InCoder-6B ficou em 0 de Pass@1 em Java com e sem contexto (S137-C1 a S137-C3).

O resultado parece depender do recuperador e da qualidade do que se recupera (S273, S547, S1195). As condições oráculo de S120 e S556 dão limite superior, não condição realista, e há recuperadores treinados em S120, S170 e S852.

\paragraph{Instrução, restrições e persona}
Avaliam-se especificação, pré e pós-condições, restrições de segurança ou acessibilidade, regras de estilo, persona, redação e língua da instrução. Nas sínteses parciais, cerca de 17 estudos do grupo só com instrução, exemplos, contexto ou restrições tratam de instrução e restrições, e 9 de formulação e língua, com contrastes também nos outros grupos.

% AFIRMACAO: A86
Em dois estudos, mudanças só de redação pioraram o resultado, e a persona não ajudou: corrigir problemas de qualidade dos prompts do HumanEval Java reduziu o pass@1 do CodeGen-2.5-7B-multi de 0,219 para 0,15 (S323-C2), cinco de seis frases instrucionais elevaram as soluções incorretas do Codex de 45 para valores entre 59 e 74 em 120, e a sexta as reduziu para 44 (S2162-C1 a S2162-C6), e a persona mudou o Pass@1 do GPT-4o de 47,5\% para 47,4\% (S560-C3).
% AFIRMACAO: A87
Restrições dirigidas a requisitos não funcionais tiveram direção variável: o prompt estruturado de segurança reduziu o RPE do GPT 4o de 0,42 para 0,13 (S1607-C1), mas, com GPT-3.5-turbo, 13 variantes instrucionais sem iteração elevaram a proporção de amostras vulneráveis de 5,74\% para valores entre 5,99\% e 7,82\% (S604-C1 a S604-C10, S604-C13 a S604-C15), as diretrizes OWASP reduziram o Pass@1 de 16,19\% para 9,52\% (S1195-C4) e a instrução WCAG elevou a inacessibilidade de 0,326 para 0,433 (S2113-C1).
% AFIRMACAO: A88
A língua da instrução também alterou o resultado: com StarCoder 3B no X-HumanEval-X, a autotradução da instrução chinesa para o inglês reduziu o pass@1 de 18,29\% para 1,83\%, enquanto a tradução externa o elevou para 20,12\% e 21,95\% (S2141-C1 a S2141-C3), e, com Claude3.5, o prompt em árabe obteve Pass@1 de 0,831 contra 0,938 em inglês (S495-C1).

Scanners aplicados a \textit{benchmarks} não demonstram segurança em produção. O ganho de pré e pós-condições em S1812-C1 é agregado sobre seis modelos e, segundo a ficha, não significativo. Em S739, o objetivo da instrução muda a contagem de vulnerabilidades ao longo de dez iterações, mas sem condição sem iteração.

\paragraph{Otimização automática do prompt}
Avaliam-se a busca iterativa ou evolutiva sobre o prompt, a compressão ou poda e a seleção automática entre estilos. As sínteses parciais registram otimização em S433, S650 e S2053, compressão ou poda em S499, S1644, S1837 e S2048 e seleção em pacotes híbridos em S570, S1853 e S1924.

% AFIRMACAO: A89
Com GPT-3.5, o Prochemy elevou o Pass@1 de 72,6\% para 76,2\% no HumanEval e de 60,5\% para 73,0\% na tradução Java$\rightarrow$Python do CodeNet (S650-C1, S650-C2), e o EPiC elevou o pass@1 do o3-mini no HumanEval+ de 0,88 para 0,945 e 0,941 (proporção) nos dois registros do método (S433-C1, S2053-C1), embora as fichas registrem empate do CoT com e sem Prochemy no LiveCodeBench e o EPiC abaixo do LATS e do LDB no MBPP+ (S650; S2053).
% AFIRMACAO: A90
A compressão da docstring reduziu o Pass@1 do DeepSeekCoder-1.3B no HumanEval de 63,42\% para 54,88\%, 53,05\% e 57,93\%, conforme o método (S499-C1 a S499-C3), enquanto a poda guiada de exemplos, no ClassEval com GPT-3.5-turbo-0613, elevou o Pass@1 de 3,8\% para 9,5\% e de 4\% para 8,6\% contra o \textit{few-shot} sem poda (S1644-C1, S2048-C1).

No ClassEval, os valores absolutos ficam abaixo de 10\%, e em S2048-C1 a intervenção usa mais exemplos que o comparador. A economia de \textit{tokens} de S499 não equivale a redução de custo medida. Os pacotes híbridos incluem seleção de modelo ou de \textit{solver}, e S1853 mede o tempo de busca; nenhum desses ganhos se atribui à seleção de prompt.

\paragraph{Estudos extraídos depois das sínteses parciais}
S1610 e S1617 entraram depois da decisão do autor sobre a data de I4 e não estão nas contagens dos agrupamentos
acima.
% AFIRMACAO: A103
Em S1610, com Qwen 2.5 Coder 7B e 81 prompts, a proporção de arquivos gerados com ao menos um alerta do Bandit caiu
de 63,0\% no zero-shot para 54,3\% com \textit{few-shot} e para 55,6\% com CoT. Nas duas condições, o contraste também
acrescenta uma diretiva de segurança, e por isso é de pacote (S1610-C1; S1610-C2).
% AFIRMACAO: A104
Em S1617, com GPT-5, Gemini 2.5 Pro e DeepSeek V3.2 no HumanEval-XL, o Pass@1 agregado foi de 0,858 com prompts em inglês e de 0,844 com prompts em espanhol (S1617-C1).

\paragraph{Lacunas}
As sínteses parciais deixam lacunas que limitam os achados de geração:
\begin{itemize}
\item \textbf{Isolamento com orçamento igual.} Só S479 e S645 igualam o orçamento do laço de \textit{feedback}; nenhuma comparação contrasta multiagente e agente único com as mesmas chamadas, nem separa raciocínio de decomposição no mesmo pacote.
\item \textbf{Custo e uso real.} Nenhuma comparação mede custo ou tempo de forma pareada, e os \textit{benchmarks} não sustentam conclusões sobre produtividade, segurança ou redução de custo.
\item \textbf{Informação oráculo.} Condições com informação da solução de referência (S120, S295-C3, S399-C3, S556, S670, S892, S902-C2, S1613-C1, S1708, S2140-C2, S2145-C6) indicam limites superiores.
\item \textbf{Modelos e domínios.} Modelos pequenos (por exemplo, S114, S479, S759 e S1977) e linguagens de hardware ou de domínio aparecem em poucos estudos, com conjuntos próprios.
\item \textbf{Comparações de uma célula.} Cada comparação cobre um modelo e um \textit{benchmark}; as inversões entre modelos ficam só como heterogeneidade.
\item \textbf{Qualidade e conferência.} A contaminação foi julgada ``atende'' em 9 dos 222 estudos das sínteses parciais (S247, S323, S397, S398, S560, S782, S902, S2054 e S2123), e nenhuma ficha passou por conferência humana. S433 e S2053 descrevem o EPiC com o mesmo modelo e \textit{benchmark}; esta revisão não verificou se são versões do mesmo trabalho, o que afetaria as contagens.
\end{itemize}

\subsubsection{Depuração e reparo}

A síntese parcial de depuração cobre 59 estudos de depuração e reparo de programas, com 126 comparações extraídas nessa tarefa. No
sentido da métrica, 82 favorecem a condição com a técnica, 22 a desfavorecem e 3 empatam. As outras 19 não têm
direção definida, por não terem sentido declarado ou os dois valores. Em 46 estudos, ao menos uma técnica focal ou
componente tem avaliação do pacote, e em 38, ao menos uma tem ingrediente isolado ou contraste pareado. Outros 8 estudos não têm
comparação numérica nessa tarefa. Os agrupamentos abaixo são desta revisão, e um estudo pode entrar em mais de um.
As contagens por agrupamento vêm das sínteses parciais. Os 3 estudos de depuração incluídos na última leva da
etapa 3 estão no fim da subseção.

\paragraph{Feedback de testes e execução}
Treze estudos devolvem ao modelo, em nova rodada, resultados de testes, de depurador, de verificadores ou contraexemplos. Só quatro isolam o feedback como ingrediente (S62, S114, S412, S613). S65 traz só rodadas cumulativas, e ao menos cinco estudos o avaliam apenas dentro de pacotes.
% AFIRMACAO: A6
No ARHE, as observações reais do depurador deram ao AutoSD patch plausível em 73\% das execuções, contra 63\% com observações previstas pelo próprio LLM (S62-C2).
% AFIRMACAO: A7
Com ChatGPT no Defects4J, o retorno das falhas de testes elevou de 36 para 62 os bugs corrigidos (S412-C3); no DS-1000, retirar o conhecimento de bug derivado dos testes reduziu o ANF de 104 para 83 (S613-C3).

Há também nulos e inversões.
% AFIRMACAO: A8
Com InCoder 1B, as estruturas de reparo P1 e P2 empataram em 0,0\% (S114-C2); em S2164, a taxa AC subiu de 72\% para 73,53\% no LeetCode, mas, segundo a heterogeneidade registrada na ficha, caiu de 31\% para 29\% no Codeforces (S2164-C2).
O número de chamadas raramente é igualado (o ThinkRepair admite até 25 sessões; S412-C1), e em S65 as localizações fornecidas vêm do gabarito, isto é, são informação oráculo.

\paragraph{Informação diagnóstica no prompt}
Vinte estudos inserem uma única vez mensagem do compilador, tipo de defeito, testes que falham, traços, relato de bug ou localização.
% AFIRMACAO: A9
Em S647, a localização perfeita, informação oráculo, elevou as correções corretas de 34 para 44 (S647-C1); trocar os exemplos por teste disparador ou por relato de bug levou-as a 59 e 72, mudança que retira e acrescenta ingredientes ao mesmo tempo (S647-C2; S647-C3).
% AFIRMACAO: A10
Com Qwen2.5-Coder:7B em 40 programas do QuixBugs, descrição, linha defeituosa ou saída dos testes elevaram a taxa de reparo de 62,5\% para 70\%, enquanto o PDG e o arquivo de testes a reduziram para 55\% e 57,5\% (S1674-C1; S1674-C2; S1674-C5; S1674-C3; S1674-C4).
A linha defeituosa, nesse estudo, vem da versão correta (oráculo).
% AFIRMACAO: A11
Em S1678, acrescentar stack trace ao prompt de localização reduziu os bugs localizados de 47 para 32, e stack trace com asserção, para 41, em contagem cuja unidade de agregação não é informada (S1678-C1; S1678-C2).
% AFIRMACAO: A12
Em S233, as mensagens de execução aumentaram a Pass Rate nos bugs de sintaxe e de referência e a reduziram nos de lógica ($-$4,3 com GPT-4 e com GPT-3.5) (S233-C1); em S1008, mais contexto não teve efeito na compilação ($\chi^2 = 0{,}21$; $p = 0{,}649$) (S1008-C1).
% AFIRMACAO: A13
Entre os acréscimos favoráveis estão os traços de execução, de 154 para 177 correções corretas (S1304-C2), e as ferramentas de busca, de 14 para 23 patches plausíveis em 100 bugs (S303-C2).
Uma leitura possível, sem isolamento, é que o efeito dependa do tipo de informação e da categoria do bug, já que informação estrutural ou volumosa aparece entre as pioras de S1674 e S1678.

\paragraph{Contexto recuperado}
Quinze estudos inserem código chamador ou chamado, conhecimento de API, design rationales, memória de reparos ou padrões recuperados; oito têm contraste de ingrediente (S275, S435, S613, S783, S979, S2046, S2119, S1943).
% AFIRMACAO: A14
Com GPT-3.5 no RepoBugs, o RLCE (variante one-shot) elevou a taxa de reparo correto de 0,2258, com contexto só da função defeituosa, para 0,5968 (S275-C1).
% AFIRMACAO: A15
Na ablação de S2046, com GPT-4o no Defects4J V1.2 e localização perfeita, a busca interna por dependências elevou as correções corretas de 85 para 116, e a busca externa de padrões, para 108 (S2046-C2; S2046-C3).
% AFIRMACAO: A16
Mais contexto, porém, não equivale a ganho: injetar de uma vez todas as camadas de conhecimento deu Pass@1 de 0,38, contra 0,5 da injeção progressiva (S821-C2), e o DeepK-S ficou em 0,7449, contra 0,7466 do ThinkRepair (S2119-C2).
Os orçamentos diferem (1 candidata no BaseChatGPT e até 300 no RelRepair; S979-C1), e vários estudos supõem localização perfeita (S315, S821, S979, S2046).

\paragraph{Exemplos}
Oito estudos comparam few-shot fixo ou seleção de pares defeito-correção.
% AFIRMACAO: A17
Há melhoras com seleção de exemplos por diagnóstico (de 0,76 para 0,82; S25-C2), com dois exemplos contra completamento zero-shot, em pacote (de 6,7\% para 20,3\%; S53-C1), e com memória episódica (de 44,7\% para 48,3\%; S783-C3).
% AFIRMACAO: A18
Três comparações de exemplos fixos não são favoráveis, com queda de 0,783 para 0,643 com DeepSeek Coder 1.3B (S1598-C1), de 205 para 200 plausíveis com GPT-4 (S2057-C1) e empate em 11 (S1565-C4); em S552, a implementação de referência e o programa mais próximo reduziram os programas corrigidos de 606 para 576 e 578 (S552-C3; S552-C4).

\paragraph{Raciocínio explícito}
No reparo, as sínteses parciais registram três comparações de ingrediente, em três estudos, e seis de pacote, em quatro estudos.
% AFIRMACAO: A19
Duas das três comparações de ingrediente pioram: o CoT teve 0,5161, contra 0,5645 da instrução Detail (S275-C2), e pedir explicação antes do refinamento deu pass@1 de 40,37\%, contra 40,64\% (S354-C2); a melhora de 36 para 57 em S412 inclui seleção de exemplos com raciocínio, e não só a instrução de raciocinar (S412-C2).
% AFIRMACAO: A20
Em pacotes, com Qwen3 Coder Next em circuitos quânticos, o Pass@1 caiu de 95\% (Structured) para 45\% com CoT e 63\% com ReAct (S1429-C1; S1429-C2), e com DeepSeek Coder 1.3B o CoT reduziu o overall\_score de 0,783 para 0,613 (S1598-C2).
Essas pioras vêm de um modelo pequeno e de um domínio específico, e em S1429 o conteúdo dos prompts não é descrito.

\paragraph{Decomposição, planejamento e agentes}
Onze estudos avaliam pacotes com análise de causas, planos, ferramentas ou vários agentes; há ablações em S810, S654, S303 e S1304.
% AFIRMACAO: A21
No T3, a taxa de reparo foi de 46,7\%, contra 37,6\% sem a etapa de plano e 37,3\% sem a análise de causas, com modelo da ablação não informado (S810-C2; S810-C3).
% AFIRMACAO: A22
Em S303, retirar a máquina de estados reduziu os patches plausíveis de 23 para 18 (S303-C3), e em S1304 retirar o debate multiagente reduziu as correções corretas de 177 para 158 (S1304-C3).
% AFIRMACAO: A23
Entre sistemas, o TraceRepair com DeepSeek-V3.2 corrigiu 392 bugs, contra 291 do REINFIX com GPT-4o, mas, sob o mesmo GPT-3.5, corrigiu 224 contra 241 segundo a heterogeneidade registrada (S1304-C1); o RepairAgent superou o ChatRepair por 164 a 162, relação que se inverte no subtotal Defects4J v1.2 (S303-C1).
Com número de chamadas e backbones distintos, os ganhos de pacote não se atribuem a um mecanismo único.

\paragraph{Instruções e formato}
Oito estudos variam templates, restrições, persona, táticas de influência ou modo de edição.
% AFIRMACAO: A24
Em S2054, o prompt guiado por taxonomia elevou a CER de 50,4\% para 57,3\% com GPT-4o-mini (S2054-C4); em S1565, um estilo alternativo de comentário reduziu as falhas corrigidas de 10 para 5 (S1565-C3); e em S1556 as táticas de influência não tiveram efeito principal no SWE-bench Verified (S1556-C2).

\paragraph{Estudos extraídos depois da síntese parcial}
% AFIRMACAO: A113
Em S1869, com GPT-4 Turbo no TFix, acrescentar ao prompt o identificador da regra violada elevou o CodeBLEU de
0,444 para 0,496, e combinações de dois tipos de informação chegaram a 0,513 e 0,528 (S1869-C1 a S1869-C3).
% AFIRMACAO: A114
Em S1715, um prompt que induz memorização levou o GPT-3.5 a gerar, no Defects4J, patches idênticos ao de referência
em 78,76\% dos patches corretos, resultado que o estudo usa como indício de memorização do \textit{benchmark}
(S1715-C1). S1624 só tem resultados descritivos.

\paragraph{Lacunas}
\begin{itemize}
  \item O feedback de execução só é isolado como ingrediente em 4 dos 13 estudos, e o efeito do número de rodadas não tem suporte numérico (S1429).
  \item O raciocínio explícito reúne poucas comparações, com sinais opostos e domínios restritos.
  \item Linguagens de hardware (S342, S393, S1132), circuitos quânticos (S1429), OCaml (S2130), acessibilidade web (S2135) e reparo de segurança (S2129, com 7 casos; S391-C2) dependem de um estudo ou de conjuntos sem repetição.
  \item Das 126 comparações da tarefa, 19 não têm direção definida; S1620 imprime só variação relativa, sem acurácias absolutas.
  \item Comparadores reaproveitados de outros artigos (S321-C1, S783-C1, S2046), orçamentos desiguais e métricas de plausibilidade limitam a comparação.
  \item A contaminação não é informada em 27 dos 59 estudos da síntese parcial, e nenhuma comparação mede produtividade, custo ou segurança em uso real.
\end{itemize}

\subsubsection{Tradução}

A síntese parcial de tradução cobre 25 estudos de tradução de código, e os 2 incluídos depois estão no fim desta
parte; nenhum deles registra também a tarefa de refatoração. As fichas que sustentam esta subseção foram extraídas por IA e ainda não foram conferidas por humano. Contando cada estudo uma vez por tipo de contrafactual, com sobreposição entre os tipos, 17 dos 25 estudos avaliam ao menos um pacote, 12 isolam algum ingrediente, 6 têm desenho cumulativo e 3 (S1010, S1220 e S1836) não têm contrafactual. Os conjuntos mais frequentes são CodeNet (6 dos 25 estudos), TransCoder/GeeksforGeeks (5) e AVATAR (5). Pelo mecanismo acrescentado ao prompt, os estudos formam seis agrupamentos:
\begin{itemize}
  \item reparo com retorno de compilação, execução ou testes (11 estudos);
  \item exemplos, em quantidade, seleção ou recuperação (6);
  \item explicação, raciocínio e representação intermediária (2);
  \item instruções, restrições e contexto orientados a requisitos (6);
  \item otimização automática de prompt (2, dos quais 1 com comparação de tradução);
  \item sistemas híbridos e estudos sem contrafactual (4).
\end{itemize}

\textbf{Reparo com retorno da execução.} Em 9 dos 11 estudos, a técnica focal é avaliada como pacote; apenas S329 e S1727 isolam ingredientes por ablação.
% AFIRMACAO: A25
Entre os pacotes, a acurácia computacional (CA) passou de 54\% para 80\% no S732 (GPT-4o, de C para Rust), a aprovação de sintaxe passou de 51,9\% para 75,0\% no S1132 e a aprovação passou de 80\% para 97\% no S2118, contra um agente sem análise de compilador, embora neste último a mediana de \textit{speedup} do pacote tenha sido 0,76$\times$, contra 0,88$\times$ da condição sem análise (S732-C1; S1132-C1; S2118-C1; S2118).
% AFIRMACAO: A26
Em sentido contrário, o reparo iterativo reduziu a taxa de aprovação de 29,71\% para 16,57\% no S2051 com CodeLlama-70b, queda que a ficha registra nos nove modelos avaliados, embora as métricas das duas condições não sejam equivalentes, e não alterou o resultado no S919 (Pass de 89,5\% e 89,2\%) nem no S1163, em que a compilação ficou no teto nas duas condições (S2051-C1; S2051; S919-C3; S1163-C1).

% AFIRMACAO: A27
No S329, com LLaMA-7B, o pacote elevou a CA de 22,74\% para 28,45\%; nas ablações, os casos de teste no prompt levaram a CA de 22,74\% para 27,43\%, e o reparo guiado por execução, acrescentado a eles, de 27,43\% para 28,45\% (S329-C1; S329-C2; S329-C3).
Uma leitura possível é que, nesse modelo e conjunto, a maior parte da diferença do pacote acompanha o acréscimo dos testes, e não o reparo. Os estudos também diferem no número de rodadas (até três no S1727, até cinco no S732 e até dez no S1132), e nenhum mede custo, embora o número de chamadas ao modelo varie entre eles.

\textbf{Exemplos.} Os seis estudos comparam \textit{few-shot}, \textit{many-shot} e seleção de exemplos por similaridade ou estrutura.
% AFIRMACAO: A28
Com exemplos, a CA passou de 61,28\% para 64,28\% no S674 (Qwen2-7B), com seleção estrutural contra aleatória e cinco exemplos nas duas condições, o BLEU passou de 12,14 para 12,58 no S1032 com cinco exemplos e a aprovação de sintaxe passou de 51,9\% para 60,0\% no S1132 com recuperação simples de exemplos (S674-C1; S1032-C1; S1132-C3).
% AFIRMACAO: A29
Em sentido contrário, o prompt com um exemplo reduziu a taxa de compilação de 0,49 para 0,43 no S964 com LLaMA-3.3-70B, e no S409 o StarCoder 15.5B não variou (delta de 0,000) com um, dois ou três exemplos (S964-C1; S409).
Esses resultados sugerem, sem isolar o fator, dependência do modelo e do critério de seleção. Além disso, S409 e S1032 usam métricas de superfície (BLEU e CodeBLEU), que não medem correção funcional.

\textbf{Explicação e raciocínio.}
% AFIRMACAO: A30
No S193, com GPT-3.5 de Python para JavaScript, as três variantes de explicação antes da tradução tiveram pass@1 de 84,8\%, 84,3\% e 83,2\%, contra 85,5\% da tradução direta, embora a introdução do estudo afirme melhora consistente (S193-C1; S193-C2; S193-C3; S193).
% AFIRMACAO: A31
No S836 (Open GPT4 8X7B, CodeNet), os prompts \textit{one-shot} com raciocínio elevaram a taxa de sucesso de 28,6\% para valores entre 32,2\% e 42,4\%, ao passo que as quatro abordagens \textit{zero-shot} em duas etapas (código, depois linguagem natural ou AST, depois linguagem-alvo) a reduziram para 10,7\%, 2,6\%, 9,0\% e 11,2\% (S836-C1; S836-C4; S836-C7; S836-C8; S836-C2; S836-C3; S836-C5; S836-C6).
Como as condições com raciocínio também acrescentam um exemplo, esse ganho é de pacote e não pode ser atribuído ao raciocínio.

\textbf{Instruções e requisitos.}
% AFIRMACAO: A32
No S304, com CodeLlama 7B e a mesma expressão regular nas duas condições, a restrição de formato elevou a taxa de extração da saída (MSR) de 6,6\% para 94,2\%, mas no Mixtral a relação se inverteu (MSR de 60,4\% contra 82,20\%; CA de 0\% contra 35,60\%) (S304-C1; S304).
A MSR mede a extração da resposta, e não sua correção.
% AFIRMACAO: A33
No S951, com GPT4omini em 67 amostras, a métrica VIR (menor é melhor) foi de 100,0\% na linha de base, 74,6\% com requisitos de segurança e 67,2\% com recuperação de casos vulneráveis, sem explicação para a linha de base de 100\% e com a funcionalidade medida em outro conjunto (S951-C1; S951-C2; S951).
Esse resultado em 67 amostras não demonstra segurança em uso real.

\textbf{Otimização de prompt e sistemas híbridos.}
% AFIRMACAO: A34
No S650, com GPT-3.5, a otimização automática de prompt elevou o Pass@1 de 60,5\% para 73,0\% de Java para Python no CodeNet, com um caso de teste por tarefa, e a ablação da iteração foi feita apenas em uma tarefa de geração (S650-C2).
O S50 inclui componentes fora do escopo, e S1220 e S1836 são descritivos.

\textbf{Lacunas.} A evidência é insuficiente nos seguintes pontos:
\begin{itemize}
  \item nenhuma comparação extraída isola o raciocínio do exemplo; S193 isola a explicação, mas com um modelo e um par de linguagens;
  \item a otimização de prompt para tradução tem 1 estudo com 1 comparação (S650);
  \item a segurança tem 1 estudo (S951), com um modelo nas comparações extraídas;
  \item a migração de versões e domínios (Java 8 para 11, Python para Spark, SQL, VHDL, Triton) só aparece em estudos sem contrafactual, cumulativos ou com conjunto interno (S1220, S1836, S1803, S1132);
  \item nenhuma comparação de tradução mede custo ou número de chamadas.
\end{itemize}

\textbf{Estudos extraídos depois da síntese parcial.}
% AFIRMACAO: A115
Em S1713, com Claude-2 na tradução de 569 programas C++ para Rust, o \textit{few-shot} elevou de 240 para 539 os
programas que compilam. Os contraexemplos não acrescentaram nada além do \textit{few-shot} (539), e o pacote inclui
reparo por regras manuais (S1713-C1 a S1713-C3).
% AFIRMACAO: A116
Em S1807, com Llama 3.1 70B traduzindo programas C do CodeNet, o \textit{zero-shot} curado mudou a taxa de
traduções que não compilam de 14,91\% para 14,79\%, e o CoT a elevou para 21,36\% (S1807-C1; S1807-C2).

\subsubsection{Refatoração}

A síntese parcial de refatoração cobre 19 estudos de refatoração, e o incluído depois está no fim desta parte. Com sobreposição entre os tipos, 12 dos 19 estudos avaliam ao menos um pacote, 9 isolam algum ingrediente, 1 tem desenho cumulativo (S2155) e 4 não têm contrafactual (S739, S1368, S1698 e S1706). No S958, a linha de base é o código humano do LeetCode, e não outra condição de prompting. Quinze dos 19 estudos usam conjuntos próprios ou projetos específicos, e a contaminação consta como não informada em 13 dos 19. Os agrupamentos são:
\begin{itemize}
  \item exemplos e recuperação de exemplos (6 estudos);
  \item especificação da transformação e conhecimento de domínio (6);
  \item raciocínio explícito (2);
  \item retorno de testes ou compilação, autorreparo e iteração (5);
  \item estudos sem contrafactual (4).
\end{itemize}

\textbf{Exemplos.}
% AFIRMACAO: A35
Com exemplos e o mesmo modelo nas duas condições, a aprovação em testes passou de 28,36\% para 34,51\% no S4 (StarCoder2-15B, \textit{one-shot}), o Pass@1 passou de 87,44\% para 91,11\% no S202 (gpt-3.5-turbo-0301) e a correção passou de 87,33\% para 90,0\% no S1822 com um exemplo recuperado (DeepSeek-V3) (S4-C2; S202-C1; S1822-C1).
% AFIRMACAO: A36
Há também efeitos colaterais ou nulos: no S202, a compilabilidade caiu de 99,99\% para 99,93\% e a complexidade ciclomática do \textit{one-shot} piorou; no S1685, as compilações passaram de 31 para 32 em 37 amostras, sem teste estatístico; e no S958 o tempo de execução foi de 0,00606 para 0,00617 ms, em comparação com o código humano (S202; S1685-C2; S958-C4).

\textbf{Especificação da transformação.} Os seis estudos nomeiam o tipo de refatoração ou injetam esquemas, \textit{changelogs}, taxonomias, regras ou definições; S466, S511, S799 e S2122 isolam ingredientes.
% AFIRMACAO: A37
No S511, com GPT-4 em 180 casos, especificar o tipo de refatoração elevou o \textit{recall} de 15,6\% para 52,2\%, acrescentar subcategoria e exemplo o elevou de 52,2\% para 66,7\%, e limitar o espaço de busca o levou de 15,6\% para 86,7\%, valor que a conclusão do estudo registra como 86,1\% (S511-C1; S511-C2; S511-C3; S511).
% AFIRMACAO: A38
No S2122, o pacote elevou o Pass@1 de 4\% para 76\% com Mistral-Large, e com LLaMA-4-Maverick retirar o esquema da API ou os \textit{changelogs} deu 56\% e 60\%, contra 76\% do pacote, mas o CrystalBLEU caiu de 0,78 para 0,66 enquanto o Pass@1 subia (S2122-C1; S2122-C2; S2122-C3; S2122).
% AFIRMACAO: A39
No S466, a redução de \textit{smells} passou de 55\% para 89\%, e no S799 os verdadeiros positivos passaram de 10 para 12 em 25 trechos, com uma invocação; o ganho da DSL de regras no S466 (de 67\% para 89\%) foi relatado só em prosa, com modelo não informado (S466-C1; S466-C2; S799-C1).
% AFIRMACAO: A40
Em sentido contrário, no S362 uma linha de base por regras (RIdiom) teve precisão maior em 8 de 9 idiomas, segundo a ficha, que não traz números da comparação (S362-C1; S362).
Uma leitura possível é que os maiores ganhos aparecem quando a linha de base é muito baixa e a tarefa depende de conhecimento de versão (S511-C2; S2122-C1).

\textbf{Raciocínio explícito.}
% AFIRMACAO: A41
O raciocínio explícito elevou a aprovação de 28,36\% para 32,22\% no S4, avaliado como pacote, e no S958 a energia foi de 1,7122 para 1,7112 mWh em comparação com o código humano, o que o próprio estudo descreve como consumo quase igual (S4-C1; S958-C1; S958).
A medição de energia em benchmark não demonstra redução de custo operacional.

\textbf{Retorno e iteração.}
% AFIRMACAO: A42
Entre os resultados favoráveis, a proporção TPP passou de 0,503 para 0,638 no S1054 (CodeQwen-1.5), com um comparador que difere também no exemplo e com narrativa que contradiz a tabela, a validade sintática passou de 86,1\% para 92,4\% no S2183 e as compilações passaram de 23 para 32 no S1685, contra o GPT com prompt não descrito (S1054-C1; S2183-C1; S1685-C1).
% AFIRMACAO: A43
Há também inversões: no S1685, a condição sem o retorno do RefactoringMiner teve mais compilações (34 contra 32 do pacote), e no S2164, com GPT-4, a segunda rodada elevou a aceitação de 72\% para 73,53\%, enquanto pedir otimização aumentou o tempo de 86,24 ms para 90,53 ms (S1685-C3; S2164-C2; S2164-C1).
% AFIRMACAO: A44
No S739, sem linha de base, após 10 iterações com GPT-4o sem intervenção humana, as vulnerabilidades contadas foram 124 na estratégia de eficiência, 158 na de funcionalidade, 67 na de melhoria genérica e 38 na de segurança (S739-C1; S739-C2; S739-C3).
No S1685, o modelo não é informado e a reprodutibilidade não é atendida; no S739, as comparações são entre estratégias, e não contra uma condição sem a técnica.

\textbf{Lacunas.} A evidência é insuficiente nos seguintes pontos:
\begin{itemize}
  \item o raciocínio isolado: os 2 estudos o avaliam como pacote ou contra código humano;
  \item testes, arquitetura, Terraform, SQL e código descompilado, com 1 estudo cada (S466, S1775, S2122, S2183 e S1822);
  \item a adoção por desenvolvedores, tratada apenas pelo S1368, que é descritivo.
\end{itemize}

\textbf{Estudo extraído depois da síntese parcial.}
% AFIRMACAO: A117
S1883 relata só um ganho médio de velocidade de duas vezes ao converter, com GPT-4o, 100 códigos seriais para
execução distribuída em dois nós. O resultado é de um sistema que combina prompting e ambiente de execução, sem
comparador (S1883-C1).

